\section{Introduction}
Imagine we need to evaluate 1 million pairs of model generations---a task whose scale makes human annotation impractical, if not impossible. Today, a commonly proposed solution is to `just ask GPT-4'\citep{mt-bench, alpacafarm}, realizing a tempting idea that large language models (LLMs) may serve as a scalable substitute for manual annotation \citep{can-llms-be-an-alternative}. However, this compelling prospect comes with a crucial caveat---LLM-based evaluation would always remain, at best, an approximation of human judgement. Without a provable guarantee of reliability, it is no surprise that the judge model has to be chosen heuristically, often times to be the strongest and the most expensive model available (\eg GPT-4). Yet, prior works show that even the strongest judge models suffer from systematic biases \citep{LLMs-are-not-fair-evaluators, judging-judges} and over-confidence \citep{confidence-elicitation}, casting doubt on the dependability of these models. This raises a fundamental question: \textit{How can we guarantee the reliability of LLM-based evaluation?}
% Today, the pervasive call of ``just ask GPT-4'' \citep{mt-bench, alpacafarm} presents one tempting solution, suggesting that LLMs may serve as a scalable substitute for manual annotation \citep{can-llms-be-an-alternative}.
% LLM-based evaluation has emerged as a compelling approach for evaluating language models---where LLMs themselves are asked for a preference between a pair of model generations \citep{mt-bench, alpacafarm}. Despite its promise, LLM-based evaluation would always remain, at best, an approximation of human judgement---thus it is imperative that we understand how and when this approximation is reliable. In the absence of a provable guarantee of reliability, it is no surprise that the judge model has to be chosen heuristically, often times to be the strongest and the most expensive model available (\eg GPT-4). Yet, prior works suggest that even the strongest judge models suffer from systematic biases \citep{LLMs-are-not-fair-evaluators, judging-judges} and over-confidence \citep{confidence-elicitation}, casting doubt on the reliability of LLM-based evaluation.
% Such a task requires a scalable method, as manually collecting human annotations would be impractical.
% since manual annotations are impractical, we may want a scalable method to simulate human preferences.
% Imagine we need to evaluate 1 million pairs of model generations without manually annotating them.

% Imagine we need to evaluate 1 million pairs of model generations. Such a task requires a scalable method, as manual human annotation would be impractical. Today, the most tempting solution is to ``just ask GPT-4'' \citep{mt-bench, alpacafarm}. Yet, this solution is ambivalent: despite its prevalence, LLM-based evaluation would always remain, at best, an approximation of human judgement. This beckons a critical question---how can we guarantee the reliability of LLM-based evaluation?  Without a provable guarantee of reliability, it is no surprise that the judge model must be chosen heuristically, often times to be the strongest and the most expensive model available (\eg GPT-4). But, prior works suggest that even the strongest judge models suffer from systematic biases \citep{LLMs-are-not-fair-evaluators, judging-judges} and over-confidence \citep{confidence-elicitation}, casting doubt on the reliability of LLM-based evaluation.

In this work, we aim to improve the reliability of LLM-based evaluation by providing a rigorous guarantee of human agreement. That is, given a user-defined risk level $\alpha$, we provide a guarantee that, for an unseen instance $x$,
\begin{equation*}
    P(\textit{LLM preference on $x$ agrees with human}\,|\,\textit{LLM evaluates $x$}) \ge 1 - \alpha.
\end{equation*}
% $P(\textit{LLM preference agrees with human}) \ge 1 - \alpha$.
To provide this guarantee, we posit that a reliable evaluation framework should not only consider the preference of a model, but also the validity of the preference---\ie how likely humans would agree with the model judgement. When a model cannot confidently evaluate a given instance, we should not rely on its evaluated result. This motivates \emph{selective evaluation}: we evaluate an instance with an LLM judge, assess the confidence that humans would agree with its evaluation, then decide whether or not to trust the evaluated result. We show that under this framework, human agreement can indeed be guaranteed---both theoretically and empirically---by choosing when to trust the model through \textit{fixed sequence testing} \citep{fixed-sequence-testing} on a small calibration set.

The practicality of selective evaluation lies not only in achieving high agreement with humans, but also in maximizing the coverage of evaluated instances without abstention---a factor that depends on the quality of confidence measure. We find that existing methods for confidence estimation (\eg predictive probability) are brittle even with the strongest judge model, as they tend to overestimate human agreement. We then propose \textbf{Simulated Annotators}, a novel method to simulate diverse annotator preferences through in-context learning and estimate confidence as an agreement ratio between the simulations. Without relying on any external supervision, \emph{Simulated Annotators} significantly improves both the calibration and failure prediction of LLM judges. As a result, selective evaluation can be done with high coverage while satisfying the prescribed human agreement level.

% One rationale behind using the strongest LLM as a judge is that it would better predict what humans prefer. 
Moreover, since our framework provides a model-agnostic guarantee of human agreement, we no longer have to rely solely on GPT-4 for evaluation. We propose \textbf{\method} (Figure \ref{fig:teaser}), where we start from a substantially cheaper LM (\eg Mistral-7B) as a judge, and escalate to a stronger model only when the previous judge is not sufficiently confident---all while guaranteeing high agreement with humans. Importantly, users do not have to manually choose when to use which judge model; given a user-specified risk tolerance, the abstention policy is automatically decided to maintain risk control.

We test our method across preference domains including summarization and real-world user-chatbot interaction, and find that \textit{\method} significantly reduces the evaluation overhead while guaranteeing high agreement. For example, our method can outperform GPT-4 by achieving over $80\%$ human agreement in ChatArena \citep{preference-dissection}, while covering $79.1\%$ of all samples, among which $88.1\%$ are evaluated by substantially cheaper Mistral-7B or GPT-3.5 instead of GPT-4. We also show that our abstention policy closely aligns with the subjectivity perceived by humans, rather than relying on shallow features such as length ratio or token overlap. Overall, our work suggests a principled approach to make LLM-based evaluation more reliable yet cost-effective, without exclusively counting on the capabilities of the most advanced LLMs as judges.

\begin{figure*}[t]
    \centering
    \vspace{-20pt}
    \includegraphics[width=.96\textwidth]{figures/overview.pdf}
    \vspace{-5pt}
    \caption{Illustration of \method. We start with a small, cost-effective model as initial judge, estimate its confidence, and escalate to a stronger model only when the previous judge is not confident. By calibrating when to trust which judge model, \textbf{our method provides a rigorous guarantee of human agreement while employing substantially cheaper judge models.}}
    \vspace{-10pt}
    \label{fig:teaser}
\end{figure*}

% We start with a small, cost-effective model as an initial judge, estimate its confidence by simulating annotators, and escalate to a stronger model when it is not sufficiently confident.

% \begin{figure*}[t]
%     \centering
%     \vspace{-15pt}
%     \includegraphics[width=.93\textwidth]{figures/overview_figure_v4.pdf}
%     \vspace{-5pt}
%     \caption{TBU Illustration of Cascaded Selective Evaluation and Simulated Annotators for confidence estimation. \faeze{some font size can be increased.}.\\\\}
%     \vspace{-15pt}
%     \label{fig:teaser}
% \end{figure*}
